\section{Methodology}
The proposed methodology proceeds through four interconnected stages, each building upon the preceding phase. The process begins with collection of publicly available demand signals, continues through data preprocessing and organization, applies forecasting models to generate predictions, and concludes with translation of numerical outputs into actionable product recommendations.

\subsection{Data Collection}
Given the inaccessibility of proprietary Notion analytics, this research relies exclusively on publicly available data sources. Search interest data is obtained from Google Trends to quantify the frequency of queries related to specific Notion features such as ``Notion AI'' and ``Notion templates.'' User review datasets are obtained from Kaggle repositories, providing access to thousands of app store reviews that reflect user satisfaction and feature preferences. Additionally, basic app metadata including ratings, download counts, and version histories are collected to supplement the primary data sources.

All collected data is timestamped and organized to ensure temporal alignment at monthly intervals. The collection infrastructure addresses common technical challenges including API rate limits, network reliability issues, and comprehensive logging to enable traceability of data provenance.

\subsection{Data Preprocessing}
Raw data requires substantial preprocessing before analysis. Duplicate entries are identified and removed to prevent distortion of frequency counts. Missing values are addressed through interpolation or forward-filling techniques using the most recent available value. Statistical outliers that deviate substantially from expected patterns are flagged and filtered. Natural language processing techniques are applied to review text to quantify feature mention frequency and compute sentiment scores indicating whether user affect is positive, negative, or neutral regarding specific features.

\subsection{Forecasting Model Development}
Rather than relying on a single forecasting method, this framework applies three complementary approaches to capture different aspects of demand patterns \cite{makridakis2018statistical}. Exponential Smoothing is particularly responsive to recent trends \cite{hyndman2008forecasting}—if interest in a feature increased substantially in the preceding period, this model captures the shift rapidly. ARIMA analysis provides perspective on longer-term patterns including seasonality, cyclical behavior, and autocorrelation structures \cite{box2015time}. Linear Regression enables incorporation of exogenous factors such as competitive product launches or general industry trends.

Each model is trained on historical data with a held-out portion reserved for validation. Model outputs are then combined using an ensemble approach where individual model weights are determined by historical prediction accuracy—models demonstrating superior predictive performance receive proportionally greater influence in the final forecast.

\subsection{Actionable Recommendation Generation}
Numerical forecasts require translation into practical guidance for product teams. The framework converts forecast trajectories into categorical recommendations. Features exhibiting increasing demand trajectories are designated for user interface improvements and increased development resource allocation. Features demonstrating stable, consistent demand are recommended for maintenance and potential complementary additions. Features showing declining interest are flagged for strategic repositioning or redesign consideration.

This translation creates a direct connection from quantitative analysis to qualitative design decisions.

The visualization component serves a critical role in this process. Comprehensive plots comparing individual feature trajectories are generated, multiple forecasting methods are overlaid to illustrate areas of consensus and divergence, and resource allocation charts directly relate growth rates to budget recommendations. Visualizations are saved at publication quality (300 DPI) for inclusion in reports and presentations, ensuring stakeholders can interpret key insights without requiring analysis of raw numerical data.

\subsection{Model Evaluation}
Model performance is evaluated using a train-validation-test split methodology. For time-series data, rolling-origin testing is additionally employed, sliding the prediction window forward and assessing model performance at each step. Standard evaluation metrics are tracked: Mean Absolute Error quantifies typical prediction deviation, Root Mean Squared Error applies greater penalty to larger errors, and Mean Absolute Percentage Error enables comparison across features with different scales.

Where computationally feasible, confidence intervals are generated to communicate prediction uncertainty. These intervals indicate the range within which true values are expected to fall with specified probability, enabling decision-makers to assess the reliability of individual forecasts.

\subsection{Reproducibility}
All processing steps are documented comprehensively. Data transformation scripts, model training notebooks, and visualization generation code are available in the associated repository. Researchers seeking to replicate this analysis or adapt the methodology for different SaaS products can follow the documented procedures to obtain equivalent results.
